% !TEX root = ../../../main.tex
% 所属：第四部分：Fourier 分析 / 第7章：Dirichlet定理
% 节：7.1 预备知识
% 正文按原文件顺序保存；由 ../chapter7.tex 调用。
\section{预备知识}

\begin{theorem}\label{thm:infinitely-many-primes}
	存在无穷多个素数。
\end{theorem}

\begin{proof}
	反设所有素数只有
	\[
	p_1,p_2,\ldots,p_n.
	\]
	令
	\[
	N=p_1p_2\cdots p_n+1.
	\]
	则任意 $p_j$ 都不能整除 $N$。另一方面，$N>1$ 必有素因子，
	而这个素因子不在上述素数表中，矛盾。
\end{proof}

任意大于 $2$ 的素数都是奇数，因而必为 $4k+1$ 型或
$4k+3$ 型。Euclid 的方法还可以证明其中一类素数的无穷性。

\begin{proposition}\label{prop:infinitely-many-primes-3-mod-4}
	形如 $4k+3$ 的素数有无穷多个。
\end{proposition}

\begin{proof}
	反设这类素数只有有限多个，记为
	\[
	p_0=3,p_1,p_2,\ldots,p_n.
	\]
	令
	\[
	N=4p_1p_2\cdots p_n+3.
	\]
	显然 $N\equiv3\pmod 4$。若 $N$ 的所有素因子都为 $4k+1$ 型，
	则 $N\equiv1\pmod 4$，矛盾。因此 $N$ 有一个形如 $4k+3$
	的素因子 $q$。

	对 $1\leq j\leq n$，有 $N\equiv3\pmod{p_j}$，故
	$p_j\nmid N$；又因为 $p_1\cdots p_n$ 不被 $3$ 整除，
	所以 $3\nmid N$。于是 $q$ 不在上述素数表中，矛盾。
\end{proof}

注意，这一构造不能直接用于 $4k+1$ 型素数，因为两个
$4k+3$ 型整数的乘积可能为 $4k+1$ 型。

\begin{definition}\label{def:riemann-zeta-real}
	对 $s>1$，定义 Riemann zeta 函数
	\[
	\zeta(s)=\sum_{n=1}^{\infty}\frac1{n^s}.
	\]
\end{definition}

取 $f(x)=x^{-s}$。由积分比较，
\[
\begin{aligned}
\sum_{n=1}^{\infty}\frac1{n^s}
&\leq 1+\sum_{n=2}^{\infty}\int_{n-1}^{n}\frac{\dd x}{x^s}\\
&=1+\int_1^{\infty}\frac{\dd x}{x^s}
=1+\frac1{s-1}.
\end{aligned}
\]
特别地，对任意 $s_0>1$，上述级数在 $[s_0,\infty)$ 上一致收敛，
所以 $\zeta$ 在 $(1,\infty)$ 上连续。

\begin{theorem}[Euler 乘积公式]\label{thm:euler-product-zeta-real}
	对 $s>1$，有
	\[
	\zeta(s)=\prod_p\frac1{1-p^{-s}},
	\]
	其中乘积取遍所有素数 $p$。
\end{theorem}

\begin{proof}
	固定正整数 $M>N$。由算术基本定理，每个不超过 $N$ 的正整数
	都可唯一分解为不超过 $N$ 的素数之积，且各素因子的指数小于
	$M$，因此
	\[
	\sum_{n=1}^{N}\frac1{n^s}
	\leq
	\prod_{p\leq N}
	\left(1+\frac1{p^s}+\cdots+\frac1{p^{Ms}}\right)
	\leq
	\prod_{p\leq N}\frac1{1-p^{-s}}.
	\]
	令 $N\to\infty$，得到
	\[
	\zeta(s)\leq\prod_p\frac1{1-p^{-s}}.
	\]

	反过来，展开有限乘积，并再次利用素因子分解的唯一性，有
	\[
	\prod_{p\leq N}
	\left(1+\frac1{p^s}+\cdots+\frac1{p^{Ms}}\right)
	\leq\sum_{n=1}^{\infty}\frac1{n^s}=\zeta(s).
	\]
	先令 $M\to\infty$，再令 $N\to\infty$，便有
	\[
	\prod_p\frac1{1-p^{-s}}\leq\zeta(s).
	\]
	两式合并即得结论。
\end{proof}

\begin{proposition}\label{prop:sum-reciprocals-primes-diverges}
	取遍所有素数 $p$，则
	\[
	\sum_p\frac1p=\infty.
	\]
\end{proposition}

\begin{proof}
	对 Euler 乘积取对数，得到
	\[
	-\sum_p\log\left(1-\frac1{p^s}\right)=\log\zeta(s),
	\qquad s>1.
	\]
	由
	\[
	\log(1+x)=x+O\bigl(\abs{x}^2\bigr),
	\qquad \abs{x}\leq\frac12,
	\]
	可得
	\[
	\log\zeta(s)
	=\sum_p\frac1{p^s}
	+O\left(\sum_p\frac1{p^{2s}}\right)
	=\sum_p\frac1{p^s}+O(1).
	\]
	这里误差项在 $s>1$ 时一致有界，因为
	\[
	\sum_p\frac1{p^{2s}}
	\leq\sum_{n=1}^{\infty}\frac1{n^2}<\infty.
	\]

	另一方面，对任意正整数 $M$，
	\[
	\liminf_{s\to1^+}\zeta(s)
	\geq\lim_{s\to1^+}\sum_{n=1}^{M}\frac1{n^s}
	=\sum_{n=1}^{M}\frac1n.
	\]
	令 $M\to\infty$，可知 $\zeta(s)\to\infty$，从而
	\[
	\sum_p\frac1{p^s}\longrightarrow\infty
	\qquad(s\to1^+).
	\]
	若 $\sum\limits_p p^{-1}$ 收敛，则由 $p^{-s}\leq p^{-1}$ 可知左边
	在所有 $s>1$ 时一致有界，矛盾。因此
	$\sum\limits_p p^{-1}$ 发散。
\end{proof}

